% claim.tex -- AI4Math claim of record (AI-generated, 2026-09-30; instantiated
% from harness/templates/claim/claim.tex per SOP 08b).
% Machine-readable theorem anchors are the "% LEAN:" lines; scripts/paper-lint.sh
% and the claim audit check them. All Lean listings are verbatim, byte-identical
% contiguous excerpts of frozen artifacts of tasks/20260930-mossad42-quad:
%   statements  : formalized/02-mod3x3-statements.lean (86 lines, sha256 da4eae94...)
%   proof bodies: final/02-mod3x3-proved.lean          (139 lines, sha256 79c19719...)
% Line ranges used: statement 38-85 (defs 43-58 and 64-73, signatures 76-77,
% 81-82, 85-86); the proof is a CONTIGUOUS EXCERPT, lines 36-139.
% Verification: audit/listing-verify.txt (audit/inspect-tex.py --verify).
%
% SCOPE FENCE (must be preserved in any redistribution): the Lean layer below
% proves structural theorems about a pure integer sequence. The identification of
% that sequence with the constrained domino-tiling count is a combinatorial
% semantics bridge that is NOT proved here and is stated only as a conjecture.
% The mod-3 object of this track is the sum of tiling counts over tilings whose
% HORIZONTAL-DOMINO COUNT is congruent to 0 modulo 3; it must not be abbreviated
% to "tilings mod 3", which is the residue of the total count and a different object.
\documentclass[11pt]{article}

\usepackage[T1]{fontenc}
\usepackage[utf8]{inputenc}
\usepackage{amsmath,amssymb,amsthm}
\usepackage{graphicx}
\usepackage{hyperref}
\usepackage{xurl}
\setlength{\emergencystretch}{3em}
\input{lstlean}
\lstset{language=lean}

% --- local rendering patch (2026-09-30) ---
% tectonic/XeTeX does not apply the listings literate table to multi-byte
% characters, so the Unicode set used by the listings of THIS note is mapped to
% math glyphs as active characters via the standard \lccode/\lowercase idiom
% (ASCII-only source, no extra packages). Restricted to symbols outside xeCJK's
% own punctuation table: 00B7, 2013, 2014, 230A and 230B also occur in the
% frozen sources but are special punctuation that xeCJK sets up at load time --
% making them active first aborts the run with "Control sequence ... already
% defined" (observed 2026-09-27); those route through xeCJK + SimSun instead.
% 2026-09-30 修订（渲染面缺陷修复）：本清单此前只覆盖初稿用到的码点，后续
% 正文修订引入 ≠ ∏ ∈ ⇒ α ↦ ⟺ ①②③④ 而清单未同步 —— XeTeX 只发
% "Missing character" 警告并把字形**静默丢弃**（三道机械门都查不出）。现按
% claims/scan_listing_unicode.py 的 listing 内非 ASCII 码点盘点补齐；出包前须
% 以该盘点器无 MISSING 且编译日志 "Missing character" 计数为 0 复核。
% Rendering only: listing source bytes are untouched, so the byte-exactness
% checks are unaffected, and the \ifdefined guard leaves the pdflatex path
% identical to the template.
\ifdefined\XeTeXversion
  {\lccode`\~="2022 \lowercase{\global\catcode`~=\active \global\def~{\ensuremath{\bullet}}}}
  {\lccode`\~="2115 \lowercase{\global\catcode`~=\active \global\def~{\ensuremath{\mathbb{N}}}}}
  {\lccode`\~="2124 \lowercase{\global\catcode`~=\active \global\def~{\ensuremath{\mathbb{Z}}}}}
  {\lccode`\~="2190 \lowercase{\global\catcode`~=\active \global\def~{\ensuremath{\leftarrow}}}}
  {\lccode`\~="2192 \lowercase{\global\catcode`~=\active \global\def~{\ensuremath{\rightarrow}}}}
  {\lccode`\~="2194 \lowercase{\global\catcode`~=\active \global\def~{\ensuremath{\leftrightarrow}}}}
  {\lccode`\~="2200 \lowercase{\global\catcode`~=\active \global\def~{\ensuremath{\forall}}}}
  {\lccode`\~="2203 \lowercase{\global\catcode`~=\active \global\def~{\ensuremath{\exists}}}}
  {\lccode`\~="2227 \lowercase{\global\catcode`~=\active \global\def~{\ensuremath{\wedge}}}}
  {\lccode`\~="2228 \lowercase{\global\catcode`~=\active \global\def~{\ensuremath{\vee}}}}
  {\lccode`\~="2261 \lowercase{\global\catcode`~=\active \global\def~{\ensuremath{\equiv}}}}
  {\lccode`\~="2264 \lowercase{\global\catcode`~=\active \global\def~{\ensuremath{\leq}}}}
  {\lccode`\~="2265 \lowercase{\global\catcode`~=\active \global\def~{\ensuremath{\geq}}}}
  {\lccode`\~="22A2 \lowercase{\global\catcode`~=\active \global\def~{\ensuremath{\vdash}}}}
  {\lccode`\~="27E8 \lowercase{\global\catcode`~=\active \global\def~{\ensuremath{\langle}}}}
  {\lccode`\~="27E9 \lowercase{\global\catcode`~=\active \global\def~{\ensuremath{\rangle}}}}
  {\lccode`\~="00AC \lowercase{\global\catcode`~=\active \global\def~{\ensuremath{\neg}}}}
  {\lccode`\~="2212 \lowercase{\global\catcode`~=\active \global\def~{\ensuremath{-}}}}
  {\lccode`\~="2223 \lowercase{\global\catcode`~=\active \global\def~{\ensuremath{\mid}}}}
\fi
\ifdefined\XeTeXversion
  \usepackage{xeCJK}
  \setCJKmainfont{SimSun}
\fi

\newtheorem{theorem}{Theorem}

\title{Mod-8 congruence and an interleaving bridge for the horizontal-domino-count\\
 residue-class constrained tiling count on the $3\times n$ board\\
 {\large (Lean 4 machine-checked claim of record)}}
\author{Aurora0134\thanks{Contact: via the GitHub account Aurora0134; ORCID
placeholder --- to be filled in by the author before any upload. No personal
information is carried in this note.}}
\date{\today}

\begin{document}
\maketitle

\begin{abstract}
We record a machine-checked result about two pure integer sequences
$a_2,b_2\colon\mathbb{N}\to\mathbb{Z}$ defined by an initial segment together with
integer-coefficient recurrences: the compressed subcolumn $b_2$ satisfies the strong
congruence $b_2(k)\equiv 1 \pmod 8$ for every $k\in\mathbb{N}$, the sequence $a_2$
vanishes at every odd index, and the two \emph{independently defined} sequences agree on
the even indices through the interleaving bridge $a_2(2k)=b_2(k)$. All statements are
verified by the Lean~4 kernel under the pinned mathlib revision, with no \texttt{sorry}
and axioms confined to \{propext, \texttt{Quot.sound}\} (the whitelist also allows
\texttt{Classical.choice}, which this track does not use). Novelty: search found no
occupying record. This note is a priority claim of record, not a full paper.
\end{abstract}

\section{Introduction}
This note records a formalized result produced by an AI4Math automated pipeline, in which
propositions and proofs are generated by a language model and adjudicated solely by the
Lean~4 kernel. The result concerns pure sequence-level objects: an initial-segment
recurrence pair $(a_2,b_2)$, a strong congruence for the compressed subcolumn, and the
interleaving identity that ties the two definitions together. The note is a priority claim
of record, deposited for timestamping; the corresponding narrative paper has been
deliberately deferred by the author (see the card), so no full exposition is attempted
here. Verification strength is stated in Section~\ref{sec:verif}: strict kernel
compilation, zero \texttt{sorry}, standard axiom whitelist. The mod-3 object that motivates
the sequence is the sum of tiling counts over tilings counted by the number of horizontal
dominoes, modulo 3; the kernel layer does not prove that identification, and the sequence
is described throughout as a pure integer object. No literature survey is
included---that expectation belongs to arXiv moderation and does not apply to a Zenodo
deposit.

\section{Statement}

We first fix notation. Let $a_2\colon\mathbb{N}\to\mathbb{Z}$ be the sequence with initial
values
\[
a_2(0),\dots,a_2(11) = 1, 0, 1, 0, 1, 0, 9, 0, 57, 0, 225, 0
\]
and even-lag recurrence
\[
a_2(n) = 6a_2(n-2) - 15a_2(n-4) + 28a_2(n-6) - 15a_2(n-8) + 6a_2(n-10) - a_2(n-12),
\]
and let $b_2\colon\mathbb{N}\to\mathbb{Z}$ be the \emph{independently} defined sequence with
initial values
\[
b_2(0),\dots,b_2(5) = 1, 1, 1, 9, 57, 225
\]
and order-6 recurrence
\[
b_2(k) = 6b_2(k-1) - 15b_2(k-2) + 28b_2(k-3) - 15b_2(k-4) + 6b_2(k-5) - b_2(k-6).
\]
Neither definition refers to tilings, dominoes or counting; both are pure integer
recursions, so the theorems below are structural theorems about sequences. The following
three facts are formalized.

% LEAN: tasks/20260930-mossad42-quad :: b2_mod8  grade=完全证明
\begin{theorem}[mod-8 congruence for the compressed subcolumn; main result]
\label{thm:t2a}
For every $k \in \mathbb{N}$, $\;b_2(k) \bmod 8 = 1$.
\end{theorem}

% LEAN: tasks/20260930-mossad42-quad :: a2_interleave  grade=完全证明
\begin{theorem}[interleaving bridge; main result]
\label{thm:t2b}
For every $k \in \mathbb{N}$, $\;a_2(2k) = b_2(k)$.
\end{theorem}

% LEAN: tasks/20260930-mossad42-quad :: a2_odd_eq_zero  grade=完全证明
\begin{theorem}[vanishing at odd indices]
\label{thm:t2c}
For every $n \in \mathbb{N}$, if $n \bmod 2 = 1$ then $a_2(n) = 0$.
\end{theorem}

\begin{lstlisting}[caption={Frozen formal statement (Lean 4, verbatim from \texttt{formalized/02-mod3x3-statements.lean}, lines 38--85), with the three proof-body placeholder lines omitted (see the note below the listing). sha256 \texttt{da4eae94}...).},label={lst:stmt}]
import Mathlib

/-- **T2 主序列 a2（order-12 仅偶滞后递推）**。
初值 a2(0..11) = 1, 0, 1, 0, 1, 0, 9, 0, 57, 0, 225, 0；
递推 a2(n) = 6·a2(n−2) − 15·a2(n−4) + 28·a2(n−6) − 15·a2(n−8) + 6·a2(n−10) − a2(n−12)。 -/
def a2 : ℕ → ℤ
  | 0 => 1
  | 1 => 0
  | 2 => 1
  | 3 => 0
  | 4 => 1
  | 5 => 0
  | 6 => 9
  | 7 => 0
  | 8 => 57
  | 9 => 0
  | 10 => 225
  | 11 => 0
  | n + 12 =>
      6 * a2 (n + 10) - 15 * a2 (n + 8) + 28 * a2 (n + 6)
        - 15 * a2 (n + 4) + 6 * a2 (n + 2) - a2 n

/-- **T2 压缩子列 b2（独立定义：6 个初值 + order-6 递推）**。
初值 b2(0..5) = 1, 1, 1, 9, 57, 225；
递推 b2(k) = 6·b2(k−1) − 15·b2(k−2) + 28·b2(k−3) − 15·b2(k−4) + 6·b2(k−5) − b2(k−6)。
（**独立定义**，非 `a2 (2k)` 的派生——后者会使递推退化为定义展开，见文件头返工记录。） -/
def b2 : ℕ → ℤ
  | 0 => 1
  | 1 => 1
  | 2 => 1
  | 3 => 9
  | 4 => 57
  | 5 => 225
  | k + 6 =>
      6 * b2 (k + 5) - 15 * b2 (k + 4) + 28 * b2 (k + 3)
        - 15 * b2 (k + 2) + 6 * b2 (k + 1) - b2 k

/-- **T2-A（主定理：强同余）**：压缩子列模 8 恒为 1。 -/
theorem b2_mod8 (k : ℕ) : b2 k % 8 = 1 := by

/-- **T2-B（交织桥）**：两个独立定义的序列在偶位处相合——`a2 (2k) = b2 k`。
这是本轨的**主桥**（两个独立递推定义之间的等式，需真正归纳，非定义展开）。 -/
theorem a2_interleave (k : ℕ) : a2 (2 * k) = b2 k := by

/-- **T2-C（零方向）**：主序列的奇数位恒为 0。 -/
theorem a2_odd_eq_zero (n : ℕ) (hn : n % 2 = 1) : a2 n = 0 := by
\end{lstlisting}

\noindent This is lines 38--85 of the frozen statement snapshot with exactly three lines
removed: the three proof-body placeholder lines (one per theorem). That single, stated
elision is the deposit convention (see the sibling claims): the frozen file's bodies are
the formalization department's placeholders, and reproducing an unfinished-proof marker in
this note would be rejected by the deposit's static lint. Everything else---doc-comments
including their internal notes, definitions, binders, signatures and the \texttt{:= by}
keyword---is verbatim and in order. \texttt{audit/inspect-tex.py} checks this mechanically:
it asserts that every remaining line occurs in the source in order, and that the only lines
missing from the block are those placeholder bodies. The real, sorry-free proofs follow in
excerpt.

\section{Proof}
The complete machine proof is 139 lines long and therefore exceeds a comfortable page; the
listing below is a \emph{contiguous excerpt}, lines 36--139 of
\texttt{final/02-mod3x3-proved.lean}, which covers the whole development: the two
definitions, the main congruence Theorem~\ref{thm:t2a} with its full proof body, the
interleaving bridge Theorem~\ref{thm:t2b} with its full proof body, and the odd-index
vanishing Theorem~\ref{thm:t2c} with its full proof body. The only material left out is lines
1--35 of the file: the header comment block (1--34), which is project metadata rather than
mathematical content, and one blank line. The full text, header included, is in the deposit
under \texttt{proofs/02-mod3x3-proved.lean}. A prose reconstruction is not an obligation of
this note: narrative proofs belong to a full paper, which has been deferred.

\begin{lstlisting}[caption={Verbatim excerpt, lines 36--139 of \texttt{final/02-mod3x3-proved.lean} (139 lines, sha256 \texttt{79c19719}...); the full text is in the deposit's \texttt{proofs/}.},label={lst:proof}]
import Mathlib

/-- **T2 主序列 a2（order-12 仅偶滞后递推）**。
初值 a2(0..11) = 1, 0, 1, 0, 1, 0, 9, 0, 57, 0, 225, 0；
递推 a2(n) = 6·a2(n−2) − 15·a2(n−4) + 28·a2(n−6) − 15·a2(n−8) + 6·a2(n−10) − a2(n−12)。 -/
def a2 : ℕ → ℤ
  | 0 => 1
  | 1 => 0
  | 2 => 1
  | 3 => 0
  | 4 => 1
  | 5 => 0
  | 6 => 9
  | 7 => 0
  | 8 => 57
  | 9 => 0
  | 10 => 225
  | 11 => 0
  | n + 12 =>
      6 * a2 (n + 10) - 15 * a2 (n + 8) + 28 * a2 (n + 6)
        - 15 * a2 (n + 4) + 6 * a2 (n + 2) - a2 n

/-- **T2 压缩子列 b2（独立定义：6 个初值 + order-6 递推）**。
初值 b2(0..5) = 1, 1, 1, 9, 57, 225；
递推 b2(k) = 6·b2(k−1) − 15·b2(k−2) + 28·b2(k−3) − 15·b2(k−4) + 6·b2(k−5) − b2(k−6)。
（**独立定义**，非 `a2 (2k)` 的派生——后者会使递推退化为定义展开，见文件头返工记录。） -/
def b2 : ℕ → ℤ
  | 0 => 1
  | 1 => 1
  | 2 => 1
  | 3 => 9
  | 4 => 57
  | 5 => 225
  | k + 6 =>
      6 * b2 (k + 5) - 15 * b2 (k + 4) + 28 * b2 (k + 3)
        - 15 * b2 (k + 2) + 6 * b2 (k + 1) - b2 k

/-- **T2-A（主定理：强同余）**：压缩子列模 8 恒为 1。 -/
theorem b2_mod8 (k : ℕ) : b2 k % 8 = 1 := by
  refine Nat.strong_induction_on k (fun k ih => ?_)
  by_cases hlt : k < 6
  · interval_cases k <;> norm_num [b2]
  · obtain ⟨j, rfl⟩ : ∃ j, k = j + 6 := ⟨k - 6, by omega⟩
    have h5 : b2 (j + 5) % 8 = 1 := ih (j + 5) (by omega)
    have h4 : b2 (j + 4) % 8 = 1 := ih (j + 4) (by omega)
    have h3 : b2 (j + 3) % 8 = 1 := ih (j + 3) (by omega)
    have h2 : b2 (j + 2) % 8 = 1 := ih (j + 2) (by omega)
    have h1 : b2 (j + 1) % 8 = 1 := ih (j + 1) (by omega)
    have h0 : b2 j % 8 = 1 := ih j (by omega)
    rw [show b2 (j + 6) = 6 * b2 (j + 5) - 15 * b2 (j + 4) + 28 * b2 (j + 3)
          - 15 * b2 (j + 2) + 6 * b2 (j + 1) - b2 j from by simp only [b2]]
    omega

/-- **T2-B（交织桥）**：两个独立定义的序列在偶位处相合——`a2 (2k) = b2 k`。
这是本轨的**主桥**（两个独立递推定义之间的等式，需真正归纳，非定义展开）。 -/
theorem a2_interleave (k : ℕ) : a2 (2 * k) = b2 k := by
  refine Nat.strong_induction_on k (fun k ih => ?_)
  by_cases hlt : k < 6
  · interval_cases k <;> decide
  · obtain ⟨j, rfl⟩ : ∃ j, k = j + 6 := ⟨k - 6, by omega⟩
    have e5 : a2 (2 * j + 10) = b2 (j + 5) := by
      rw [show 2 * j + 10 = 2 * (j + 5) by ring]
      exact ih (j + 5) (by omega)
    have e4 : a2 (2 * j + 8) = b2 (j + 4) := by
      rw [show 2 * j + 8 = 2 * (j + 4) by ring]
      exact ih (j + 4) (by omega)
    have e3 : a2 (2 * j + 6) = b2 (j + 3) := by
      rw [show 2 * j + 6 = 2 * (j + 3) by ring]
      exact ih (j + 3) (by omega)
    have e2 : a2 (2 * j + 4) = b2 (j + 2) := by
      rw [show 2 * j + 4 = 2 * (j + 2) by ring]
      exact ih (j + 2) (by omega)
    have e1 : a2 (2 * j + 2) = b2 (j + 1) := by
      rw [show 2 * j + 2 = 2 * (j + 1) by ring]
      exact ih (j + 1) (by omega)
    have e0 : a2 (2 * j) = b2 j := ih j (by omega)
    rw [show 2 * (j + 6) = 2 * j + 12 by ring]
    rw [show a2 (2 * j + 12) = 6 * a2 (2 * j + 10) - 15 * a2 (2 * j + 8) + 28 * a2 (2 * j + 6)
          - 15 * a2 (2 * j + 4) + 6 * a2 (2 * j + 2) - a2 (2 * j) from by simp only [a2]]
    rw [e5, e4, e3, e2, e1, e0]
    simp only [b2]

/-- **T2-C（零方向）**：主序列的奇数位恒为 0。 -/
theorem a2_odd_eq_zero (n : ℕ) (hn : n % 2 = 1) : a2 n = 0 := by
  revert hn
  refine Nat.strong_induction_on n (fun n ih hn => ?_)
  by_cases hlt : n < 12
  · interval_cases n <;> simp_all [a2]
  · obtain ⟨m, rfl⟩ : ∃ m, n = m + 12 := ⟨n - 12, by omega⟩
    have h1 : (m + 10) % 2 = 1 := by omega
    have h2 : (m + 8) % 2 = 1 := by omega
    have h3 : (m + 6) % 2 = 1 := by omega
    have h4 : (m + 4) % 2 = 1 := by omega
    have h5 : (m + 2) % 2 = 1 := by omega
    have h6 : m % 2 = 1 := by omega
    have e1 : a2 (m + 10) = 0 := ih (m + 10) (by omega) h1
    have e2 : a2 (m + 8) = 0 := ih (m + 8) (by omega) h2
    have e3 : a2 (m + 6) = 0 := ih (m + 6) (by omega) h3
    have e4 : a2 (m + 4) = 0 := ih (m + 4) (by omega) h4
    have e5 : a2 (m + 2) = 0 := ih (m + 2) (by omega) h5
    have e6 : a2 m = 0 := ih m (by omega) h6
    simp only [a2]
    rw [e1, e2, e3, e4, e5, e6]
    ring
\end{lstlisting}

\section{Verification evidence}
\label{sec:verif}
\begin{itemize}
\item Toolchain: Lean 4 \texttt{v4.34.0}; mathlib4 \texttt{v4.34.0}, revision \texttt{5ed29652}.
\item Statement integrity: \texttt{formalized/02-mod3x3-statements.lean}, 86 lines,
sha256 \texttt{da4eae94971f8a5e8f1d74f973183b1643158bb7ff96e88db97a09832db50901}.
Final artifact \texttt{final/02-mod3x3-proved.lean}, 139 lines,
sha256 \texttt{79c19719835ac4ccbffcc72707fd81dde2e52d32b8e55a4fe89dad1341c70901}.
Every listing above is byte-identical to its source (checked by \texttt{audit/inspect-tex.py},
which asserts that each listed line occurs in its source in order and that the only omitted
lines are the statement file's placeholder proof bodies).
\item Kernel check: strict compilation exits 0 with no \texttt{sorry} and no \texttt{admit}.
\texttt{\#print axioms} output, verbatim:
\begin{verbatim}
'b2_mod8' depends on axioms: [propext, Quot.sound]
'a2_interleave' depends on axioms: [propext, Quot.sound]
'a2_odd_eq_zero' depends on axioms: [propext, Quot.sound]
\end{verbatim}
All three are subsets of \{propext, \texttt{Quot.sound}, \texttt{Classical.choice}\}; this
track does not use \texttt{Classical.choice} at all, and \texttt{sorryAx} does not occur.
\item Audit chain: statement freeze (sha256), byte-wise symmetric-difference comparison against
the frozen snapshot, strict-mode compilation, and the axiom whitelist check. Originals are in
the deposit under \texttt{audit/}.
\item Gate-2 rework (recorded because it changed the theorem set). The first frozen version
defined \texttt{b2} as the derived subcolumn \texttt{a2 (2 * k)}. Because \texttt{a2} itself
is a recurrence with even lags only, the then-listed theorem asserting the order-6
recurrence of \texttt{b2} reduced to a definitional unfolding, and the gate-2 bare probes
\texttt{rfl}, \texttt{trivial} and \texttt{aesop} each closed it in one tactic: the pipeline's
degeneracy criterion flagged it as a degenerate proposition with no mathematical content.
The rework made \texttt{b2} an independent definition (six initial values plus its own
order-6 recurrence), withdrew the derived recurrence theorem---it is now the definition of
\texttt{b2}, not a theorem---and added Theorem~\ref{thm:t2b}, the interleaving bridge between
two independent definitions, which needs genuine induction and is not a one-line
\texttt{simp}. After the rework the bare probes return zero passes.
\end{itemize}

\section{Novelty evidence}
Under the pipeline's C9 novelty discipline the verdict for this object is \emph{no occupying
record found}. The full query log, including every zero-hit query, is in the deposit at
\texttt{audit/c9-record.md}; it is reproduced verbatim from the screening batch and the
present batch's supplementary search.

Summary of channels. OEIS: the sequence, its compressed subcolumn, its even-index variant,
its first difference and its partial sums all returned zero hits. The odd-index variant
matched ten entries, but every one of them is a zero sequence or a characteristic function
(all-zero strings of the kind the database returns for identically zero input), which is a
trivial artifact and not an occupying record. Channel availability was established by
positive controls before any zero-hit was read: A001835 and A005178 (the unconstrained
$3\times n$ and $4\times n$ domino-tiling counts) matched, as did A000045. Literature: four
channels were queried (arXiv API, Crossref, Mathematics Stack Exchange, zbMATH Open); there
were no direct hits on this constrained-selection object. Library layer: no hits in mathlib,
compfiles, or the scanned external repositories.

Two honest limitations are recorded and must not be dropped in any restatement.
First, the OpenAlex channel was \emph{unreachable} from the machine used for this search
(Cloudflare-fronted hosts timed out; DNS resolution was verified correct against a public
resolver, so this is not a DNS-poisoning artifact). Crossref was used as a substitute
DOI-level index. Coverage is therefore \emph{lower} than in the screening batch, and this is
a channel gap, not strengthened evidence.
Second, a zero hit is \emph{not} novelty: the strongest claim licensed here is that no
record of this constrained-selection object was found in the channels queried.
No moment or expectation novelty is claimed: mixed moments $E[V^a H^b]$ for general $m\times n$
boards are already occupied by \emph{Statistics of Domino Tilings on a Rectangular Board},
Fibonacci Quarterly (2019), doi:10.1080/00150517.2019.12427625~\cite{fq2019}.
The $3\times n$ and $4\times n$ objects of this family are $m$-dimensional variants of one
another and are not used as independent novelty evidence for each other; the variant partner
of the present $3\times n$ mod-3 track is the $4\times n$ mod-3 track, which is not cited
here as support.

\section{Scope and limitations}
The kernel-level results concern the pure integer sequences $a_2$ and $b_2$ defined above.
They do \emph{not} establish that $a_2(n)$ equals the sum of the domino-tiling counts of the
$3\times n$ board over those tilings counted by the number of horizontal dominoes, modulo 3.
That identification is a combinatorial semantics bridge; it is stated here only as a
conjecture, is not part of the theorem layer, and is supported only by numerical probes and
external sequence-database evidence. Accordingly, the results are described throughout as
structural theorems about sequences, never as a proved statement about tiling counts.
The mod-3 object is exactly the sum over tilings whose horizontal-domino count is congruent
to $0$ modulo $3$; it is not the residue of the total tiling count modulo $3$, which is a
different object and is not what is formalized here.

Two points about the theorem set must be stated plainly. First, the withdrawn item: the
earlier derived recurrence for $b_2$ was removed because it was degenerate, so this note does
not claim a proved order-6 recurrence; the order-6 recurrence is now simply the definition of
$b_2$, and the three recorded theorems are the entire content of the track. Second, the
rework history recorded in Section~\ref{sec:verif}: the first frozen version defined $b_2$
from $a_2$, which made its recurrence a mere definitional unfolding; that version failed the
degeneracy gate and was replaced by an independent definition plus a genuine interleaving
bridge between two independent definitions.

The grading of all three theorems is \emph{complete proof} (0 \texttt{sorry}, axiom
whitelist), and no weakening or extra hypothesis was used. This note contains no literature
survey; it is a deposit record, not a survey paper.

\section*{Data availability}
Lean sources, verification and novelty-search evidence are deposited on Zenodo
(DOI: \texttt{RESERVED-DOI}) and mirrored as a public snapshot repository. The note is
released under CC BY 4.0 and the code under MIT.

\section*{Use of generative AI}
\begin{itemize}
\item \emph{AI-performed stages}: proposition generation, natural-language-to-Lean
formalization, proof search, and drafting of this text.
\item \emph{Model, node, budget}: pipeline node \texttt{anthropic/a6api-main/deepseek-v4.1-flash}
(wire-id; resolved at level L2 of the pipeline's model-detection ladder, snapshot taken
2026-09-30 09:35). Source-task budget: sampling 1 of 8 allowed runs on this track (10 of 32
across four parallel tracks); judgement-class LLM calls 6 of 16 across the batch; Lean
formal compilations 4 of 12 on this track (19 of 48 across the batch); LaTeX 0. The numbers are
copied from the source task's \texttt{budget.log}.
\item \emph{Final adjudication}: all statements and proofs reported here were checked by the
Lean~4 kernel. Every proof compiles with no \texttt{sorry}, and \texttt{\#print axioms}
reports only \{propext, \texttt{Quot.sound}\} for the three theorems. Statement text was
frozen before proving and compared byte-wise afterwards.
\item The author reviewed the material and is responsible for all content. AI systems are
not listed as authors.
\end{itemize}

\section*{Author contributions}
The author determined the research question and the scope fences, verified the semantics of
the statements, and is responsible for signing off and for any deposit or release. The
automated pipeline (LLM-driven) generated the proposition, the formalization, the proof
search and the draft text. No external release is performed by the pipeline.

\begin{thebibliography}{9}
% Only works actually cited; each was checked to exist before entry. Verification
% criterion is the returned title (not merely the HTTP status), following the
% 2026-09-28 erratum lesson where a DOI tail was mis-resolved to a different article.
\bibitem{lean4}
L.~de Moura and S.~Ullrich.
\newblock The Lean 4 theorem prover and programming language.
\newblock \emph{CADE-28}, 2021. \url{https://doi.org/10.1007/978-3-030-79876-5_37}

\bibitem{mathlib}
The mathlib Community.
\newblock The Lean mathematical library.
\newblock \emph{CPP 2020}, 2020. \url{https://doi.org/10.1145/3372885.3373824}

\bibitem{fq2019}
\emph{Statistics of Domino Tilings on a Rectangular Board}.
\newblock Fibonacci Quarterly, 2019.
\newblock \url{https://doi.org/10.1080/00150517.2019.12427625}
\end{thebibliography}

\end{document}
